\section{Exact identity certificates}\label{sec:certificates}

The ancillary directory contains the five scripts cited in this note,
together with a validator that compares the certificates printed here
with their output. They use
Python~3, its standard library, and SymPy; all numerical tests are
exact rational or integer calculations. The README gives serial
commands with the resource bounds. In particular:
\begin{itemize}
\item \ancfile{identities.py} checks the homogeneous determinant
and circle-kernel identities and the barycentric moment formulas.
\item \ancfile{structure.py} checks the $64$ seven-node sign spaces,
their projection dimensions, and an algebraic-diagonal certificate.
\item \ancfile{eight_factory.py} checks the $128$ eight-node sign
spaces and the universal two-group graph identity.
\item \ancfile{second_jet_certificate.py} checks
\eqref{eq:appendix-determinant} by two determinant algorithms and
saves an exact left-null-vector certificate.
\item \ancfile{elliptic.py} checks the polynomial identities
\eqref{eq:curve-identity}, \eqref{eq:invariant-ideal}, and
\eqref{eq:biquadratic}; the rational orbit; and every cross-ratio
difference in \eqref{eq:cross-signs}. It also performs $21$ rational
tests of the simultaneous invariant.
\end{itemize}
Rational expression tests prove the specified nonvanishing and
nonconstancy claims. The general statements about all isotropic
spaces are proved in Section~\ref{sec:subspaces}, and the invariant
is proved by its polynomial identity, rather than by the $21$
sample tests.

\begingroup
\interdisplaylinepenalty=10000
For an explicit certificate of \eqref{eq:invariant-ideal}, the
two polynomial quotients are
\begin{align*}
Q_1=-20(u-3)^2\bigl(&5u^4x-15u^4+5u^3x^2-45u^3
 -54u^2vx+10u^2vy+112u^2v\\
&-45u^2x^2+135u^2x+54uvx^2-100uvx-60uvy
 +114uv\\
&-40ux+120u-20v^2x^2+120v^2x-180v^2
 -12vx^2\\
&-114vx+90vy+20x^2-180\bigr),
\end{align*}
and
\begin{align*}
Q_2=-20(x-3)^2\bigl(&-15u^2x^3-45u^2x^2+54u^2xy
 -12u^2y+60u^2\\
&+5ux^4+120ux^3-54ux^2y+135ux^2
 -100uxy-40ux\\
&-114uy-240u+10vx^2y-60vxy+90vy
 -15x^4\\
&-225x^3+112x^2y+114xy+120x+180\bigr).
\end{align*}
\endgroup
Substitution into \eqref{eq:invariant-ideal} and expansion gives
an identity in $\Q[x,u,y,v]$. The script obtains the same quotients
by successive polynomial division in $y$ and $v$ and verifies the
zero remainder.
